% generator_I06.tex -- Prop I.6 (Theor.): equal angles => equal opposite sides.
% Geometry and text from jemmybutton (CC-BY-SA). Build: lualatex generator_I06.tex
\documentclass[12pt]{article}
\usepackage{geometry}\geometry{a5paper, margin=1.5cm}
\usepackage[lining]{ebgaramond}
\usepackage{amssymb}
\usepackage{byrne}
\def\mpPre{u := 1cm; textLabels := false;}

\begin{document}

\defineNewPicture[1/4]{%
  pair A, B, C, D;%
  A := (0, 0);%
  B := A shifted (7/2u, 0);%
  D := 1/2[A,B] shifted (0, 3u);%
  C := 2/3[A, D];%
  byAngleDefine(B, A, D, byyellow, SOLID_SECTOR);%
  byAngleDefine(A, B, D, byblack,  SOLID_SECTOR);%
  draw byNamedAngleResized();%
  byLineDefine(B, C, byyellow, SOLID_LINE, REGULAR_WIDTH);%
  byLineDefine(A, B, byred,    SOLID_LINE, REGULAR_WIDTH);%
  byLineDefine(B, D, byblue,   SOLID_LINE, REGULAR_WIDTH);%
  byLineDefine(C, A, byblack,  SOLID_LINE, REGULAR_WIDTH);%
  byLineDefine(C, D, byblack,  DASHED_LINE, REGULAR_WIDTH);%
  draw byNamedLine(BC);%
  draw byNamedLineSeq(0)(CA,CD,BD,AB);%
  draw byLabelsOnPolygon(A, C, D, B)(ALL_LABELS, 0);%
}

% STATEMENT
\noindent
\begin{minipage}[t]{0.45\linewidth}
\centering\vspace{0pt}%
\drawCurrentPicture
\end{minipage}\hfill
\begin{minipage}[t]{0.52\linewidth}
\vspace{0pt}%
\textbf{Book~I.\enspace Prop.\ VI. Theor.}

\medskip\noindent\itshape
In any triangle (\drawLine[bottom][triangleABD]{CA,CD,BD,AB})
if two angles (\drawAngle{A} and \drawAngle{B}) are equal,
the sides (\drawUnitLine{CA,CD} and \drawUnitLine{BD})
opposite to them are also equal.

\bigskip

% PROOF
\noindent\upshape
For if the sides be not equal, let one of them
\drawUnitLine{CA,CD} be greater than the other \drawUnitLine{BD},
and from it cut off $\drawUnitLine{CA} = \drawUnitLine{BD}$~(pr.~I.3),
draw \drawUnitLine{BC}.

\medskip\begin{center}
Then in \drawLine[bottom]{BC,AB,CA} and \triangleABD,\\
$\drawUnitLine{CA} = \drawUnitLine{BD}$~(const.),\\
$\drawAngle{A} = \drawAngle{B}$~(hyp.)\\
and \drawUnitLine{AB} common,\\
$\therefore$ the triangles are equal~(pr.~I.4)\\
a part equal to the whole, which is absurd;\\[4pt]
$\therefore$ neither of the sides \drawUnitLine{CA,CD}
or \drawUnitLine{BD} is greater than the other,\\
hence they are equal.
\end{center}

\begin{flushright}Q.\ E.\ D.\end{flushright}
\end{minipage}

\end{document}
